% Introduction. Written before the main experiments: it states the question, the design and
% what the paper provides, and claims no model result. tests/test_paper_intro.py checks that
% it cites only bibliography keys, quotes no measured number, and makes no result claim.
% The final paragraph is to be rewritten once M5/M6 results exist.
\section{Introduction}
\label{sec:intro}

Vision-language models (VLMs) answer many questions about natural images, yet questions
about drawn structure, such as whether one node of a diagram can reach another, remain hard
for them~\cite{mcdgraph,graphverse}. Probing suggests the needed information is often
present in their representations even when the answer is wrong~\cite{signals}. A common
remedy is supervised fine-tuning on the questions themselves or on structural auxiliary
tasks such as describing or completing the graph~\cite{mcdgraph}.

We ask whether a different signal helps: \emph{learning to repair} a drawing. A repair task
gives the model a drawn directed graph and reachability goals and asks for the smallest
edge edit that satisfies them. Producing a correct edit requires reading which edges exist
and which paths depend on which edge; unlike a yes/no answer, the edit can be checked
exactly and leaves little room for guessing. If repair training builds a usable internal
reading of the graph, it should transfer to questions the model was never trained to
answer, such as reachability and edge criticality.

Testing this claim fairly needs two things that the graph-image benchmarks we reviewed
(Section~\ref{sec:related}) do not provide together. First, labels and repair answers must be exact, and the verifier must
accept every optimal edit, not one reference answer. Second, the evaluation must rule out
answers that ignore the image. We build a procedural benchmark with an exact solver, an
independent reference check and \emph{answer-flipping twins}: every question is asked on
two drawings that differ by one edge and have opposite answers, so a predictor that does
not use the image scores at most $0.50$ per task, and answering both twins correctly shows
that the deciding edge was used (Fig.~\ref{fig:twins}).

On this benchmark we compare, on the same base model with LoRA~\cite{lora} and matched
samples, root graphs and images, a repair-trained arm against question-answer training on
the target questions and against structural fine-tuning in the style of
MCDGraph~\cite{mcdgraph}. Oracle conditions that give the true graph as text, and a
parse-then-solve arm in the spirit of describe- or recover-then-compute
methods~\cite{visiongraph,vgcompiler}, help separate failures of reading the drawing from failures of
reasoning about it, within limits we state (the text oracle also changes the input format). All splits are grouped by root graph with no isomorphic graph shared across
splits, and the test splits, including held-out styles, families and sizes, are opened once.

\paragraph{Contributions.}
(i)~A verified generator and benchmark for drawn directed graphs with exact labels,
answer-flipping twins, invariance pairs and leakage-free splits (Section~\ref{sec:benchmark}).
(ii)~A pre-specified, exposure-matched protocol for testing whether edit training transfers
to non-edit questions in a small VLM (Sections~\ref{sec:setup} and~\ref{sec:analysis}).
(iii)~Oracle-graph and parse-and-solve conditions that help separate perception errors from
reasoning errors, with their limits stated. (iv)~A pre-registered external check on drawings by other authors
(VGCure) and an item-level account of why it and our locked test disagree about repair
training (Sections~\ref{sec:external} and~\ref{sec:why}). (v)~A record of what did not work
during development (Section~\ref{sec:negative}). Code and data: \url{https://github.com/efeerdogmus0/trace-graph-repair}.
% Main findings paragraph: scripts/paper_framing.py (rule framing-v2, chosen after reading); the pre-registered
% finding-v1 paragraph (scripts/paper_finding.py) is kept in reports/final/intro_finding.tex.
\IfFileExists{../reports/final/framing_intro.tex}{% generated by scripts/paper_framing.py, rule framing-v2 (chosen after reading, see DECISIONS 2026-11-23T20:30Z); do not edit
\paragraph{Main findings.} On an adapted $500$-graph VGCure evaluation (drawings by other authors in a style none of our adapters saw, with our prompts, images scaled to our input size and strict scoring), the repair-trained adapter (arm~D) answers edge-existence and reachability questions better than direct question-answer fine-tuning (arm~B): the difference is $+0.056$ (95\% interval [$+0.038$, $+0.075$], mean of 3 seeds, above zero at every seed). The lead comes from both tasks: edge existence $+0.049$ and reachability $+0.063$. On our own pre-registered locked test, which decides the hypothesis, D does worse than B: D$-$B is $-0.105$ (95\% interval [$-0.117$, $-0.093$]), so the hypothesis is not supported at this scale. Section~\ref{sec:why} traces the disagreement: on our test D's deficit is in the critical-edge task, where it answers \emph{true} to 79\% of the questions, while its reachability is level with B's; on VGCure, which has no critical-edge items, D gains on false edge claims and unreachable pairs. This is a pilot-scale study: all adapters were trained on $512$ synthetic root graphs with one 2B base model (Qwen3.5-2B, LoRA), and nothing here measures real diagrams. Both evaluations were fixed before any of their outputs existed; giving the external result this place in the abstract and introduction was decided after both were read, while the wording rule fixed in advance put only the locked-test statistic here (Section~\ref{sec:why}). Per-split, per-task and secondary comparisons are in Section~\ref{sec:test-results}.
}{%
\IfFileExists{../reports/final/intro_finding.tex}{% generated by scripts/paper_finding.py, rule finding-v1, from reports/final/final_primary.json; do not edit
\paragraph{Main finding.} Repair training (arm~D) did worse than the reference fine-tuning arm on the pre-registered six-cell statistic; the hypothesis is not supported at this scale. On the locked test the difference D$-$B is $-0.105$ (95\% interval [$-0.117$, $-0.093$], mean of 3 seeds). This is a pilot-scale study: all adapters were trained on $512$ synthetic root graphs with one 2B base model (Qwen3.5-2B, LoRA), and nothing here measures real diagrams. Per-split, per-task and secondary comparisons are in Section~\ref{sec:test-results}.
}{}}
